% Irodov Part IV "Oscillations and Waves", §4.1, independent problems 1--20.
% Original uploaded PDF pp.162--164, printed pp.153--155;
% answers PDF pp.322--323, original answers pp.313--314.
\oscproblem{1}{153}
质点沿 $x$ 轴作简谐振动，
$x=a\cos(\omega t-\pi/4)$。
作出(a)位移、速度、加速度随时间变化的示意图；
(b)速度和加速度随坐标 $x$ 变化的图。
\exsolution
求导得
\[
v_x=-a\omega\sin(\omega t-\pi/4),\qquad
a_x=-a\omega^2\cos(\omega t-\pi/4)=-\omega^2x.
\]
三条时间曲线同频率；$v_x$ 比 $x$ 的相位超前 $\pi/2$，$a_x$ 与 $x$ 反相。
消去相位：
\[
\boxed{\left(\frac{x}{a}\right)^2+
\left(\frac{v_x}{a\omega}\right)^2=1},
\qquad
\boxed{a_x=-\omega^2x}.
\]
因此 $(x,v_x)$ 为半轴 $a,a\omega$ 的椭圆，$(x,a_x)$ 为负斜率直线。
\excheck $x=\pm a$ 处 $v_x=0$；$x=0$ 处速度模最大、加速度为零，均与曲线吻合。

\oscproblem{2}{153}
质点沿 $x$ 轴运动，
$x=a\sin^2(\omega t-\pi/4)$。
求(a)振动幅度、周期与 $x(t)$ 图形；
(b)速度 $v_x$ 与坐标 $x$ 的关系及图形。
\exsolution
利用 $\sin^2u=(1-\cos2u)/2$：
\[
x(t)=\frac a2\left[1-\cos(2\omega t-\pi/2)\right],
\]
故它围绕平衡位置 $x_c=a/2$ 振动，
\[
\boxed{A=\frac a2,\qquad T=\frac\pi\omega}.
\]
速度
$v_x=a\omega\sin(2\omega t-\pi/2)$。
因为
$\sin^22u=4\sin^2u(1-\sin^2u)$，
\[
\boxed{v_x^2=4\omega^2x(a-x)},\qquad
\boxed{v_x=\pm2\omega\sqrt{x(a-x)}}.
\]
\excheck 横坐标只能为 $0\le x\le a$；相图闭合在 $x=0,a$，速度模在那里为零。

\oscproblem{3}{154}
质点沿 $x$ 轴作简谐振动，角频率
$\omega=4.00\,\mathrm{s^{-1}}$；
某时刻位置 $x_0=25.0\,\mathrm{cm}$、
速度 $v_0=100\,\mathrm{cm/s}$。
求 $t=2.40\,\mathrm s$ 后的位置与速度。
\exsolution
由运动方程 $\ddot x=-\omega^2x$，
以该时刻为 $t=0$ 的初值解为
\[
x(t)=x_0\cos\omega t+\frac{v_0}{\omega}\sin\omega t,
\quad
v(t)=v_0\cos\omega t-\omega x_0\sin\omega t.
\]
代 $\omega t=9.60$ rad、$x_0=0.25$m、$v_0=1.00$m/s：
\[
\boxed{x(2.40)\simeq-0.29\,\mathrm m},\quad
\boxed{v(2.40)\simeq-0.81\,\mathrm{m/s}}.
\]
\excheck 能量式
$v^2+\omega^2x^2=v_0^2+\omega^2x_0^2=2.0\,\mathrm{m^2/s^2}$
在前后两时刻相等。

\oscproblem{4}{154}
某质点作简谐振动，在距平衡位置分别为
$x_1,x_2$ 时的速度模为 $v_1,v_2$。
求角频率 $\omega$ 和振幅 $a$。
\exsolution
简谐振动能量关系
$v_i^2+\omega^2x_i^2=\omega^2a^2$。
两式相减得到
\[
\boxed{\omega=
\sqrt{\frac{v_1^2-v_2^2}{x_2^2-x_1^2}}}.
\]
代回任一式：
\[
\boxed{a=
\sqrt{\frac{v_1^2x_2^2-v_2^2x_1^2}
{v_1^2-v_2^2}}}.
\]
\excheck 只有上述根号内为正且数据物理相容时才有实数频率和振幅；$a\ge|x_1|,|x_2|$。

\oscproblem{5}{154}
质点简谐振动周期 $T=0.60$s、振幅
$a=10.0$cm。从(a)极端位置；
(b)平衡位置出发，
在走过路程 $a/2$ 的时间内求平均速率。
\exsolution
(a) 从极端位置设 $x=a\cos\omega t$，
走过 $a/2$ 即 $x=a/2$，
于是 $\omega t=\pi/3$，
$\Delta t=T/6$。
平均速率
$\bar v=(a/2)/(T/6)=3a/T$。
(b) 从平衡点设 $x=a\sin\omega t$，
走过 $a/2$ 即 $\sin\omega t=1/2$，
$\Delta t=T/12$，
平均速率 $6a/T$。
故
\[
\boxed{\bar v_a=\frac{3a}{T}=0.50\,\mathrm{m/s}},
\quad
\boxed{\bar v_b=\frac{6a}{T}=1.0\,\mathrm{m/s}}.
\]
\excheck 此区间内运动未反向，因此平均速率也等于位移模除以时间。

\oscproblem{6}{154}
质点从 $t=0$ 开始按
$x=a\sin\omega t$ 运动。
求起始 $3T/8$ 时间内
(a)平均速度的 $x$ 投影；
(b)平均速度矢量的模；
(c)速率的平均值。
\exsolution
$T=2\pi/\omega$，
$\tau=3T/8=3\pi/(4\omega)$，
末坐标
$x(\tau)=a\sin(3\pi/4)=a/\sqrt2$。
故
\[
\boxed{\langle v_x\rangle=
\frac{x(\tau)-x(0)}{\tau}
=\frac{2\sqrt2}{3\pi}a\omega},
\quad
\boxed{|\langle\bm v\rangle|=
\frac{2\sqrt2}{3\pi}a\omega}.
\]
这段时间先从 $0$ 到 $a$，再由 $a$ 折返至 $a/\sqrt2$，
全路程 $s=a+(a-a/\sqrt2)$，
所以
\[
\boxed{\langle|\bm v|\rangle=
\frac{s}{\tau}=
\frac{2(4-\sqrt2)}{3\pi}a\omega}.
\]
\excheck 平均速度与平均速率不同，后者计入折返路程。

\oscproblem{7}{154}
质点按 $x=a\cos\omega t$ 运动，从 $t=0$ 起经过任意时间 $t\ge0$，求累计路程 $s(t)$。
\exsolution
每经历相位 $\pi/2$，质点正好走过距离 $a$。
写 $u=\omega t$，
$n=\lfloor2u/\pi\rfloor$，
$\theta=u-n\pi/2\in[0,\pi/2)$。
完整的 $n$ 个四分之一周期累计路程为 $na$。
剩余相位中，当 $n$ 为偶数，
从端点出发，路程增量 $a(1-\cos\theta)$；
当 $n$ 为奇数，从平衡位置出发，路程增量 $a\sin\theta$。
因此
\[
\boxed{s(t)=
\begin{cases}
a[n+1-\cos\theta],&n\ \text{为偶数},\\
a[n+\sin\theta],&n\ \text{为奇数}.
\end{cases}}
\]
\excheck $t=0$ 时 $n=\theta=0$，得 $s=0$；
每完整周期路程为 $4a$。

\oscproblem{8}{154}
质点自 $t=0$ 起速度投影为
$v_x=35\cos(\pi t)\,\mathrm{cm/s}$，
其中 $t$ 以秒计。
求前 $2.80$s 内走过的路程。
\exsolution
所求路程为绝对值积分
$s=35\int_0^{2.8}|\cos\pi t|\,\mathrm dt$ cm。
速度换向时刻为 $t=0.5,1.5,2.5$s，
逐段积分，得到
\[
s=\frac{35}{\pi}
\bigl[(1)+(2)+(2)+(1-\sin0.8\pi)\bigr]
\,\mathrm{cm}.
\]
即
\[
\boxed{s=\frac{35}{\pi}(6-\sin0.8\pi)
\,\mathrm{cm}\simeq60\,\mathrm{cm}=0.60\,\mathrm m}.
\]
\excheck 不可用最终位移直接代替路程；在 $0.5,1.5,2.5$s 速度依次换向。

\oscproblem{9}{154}
质点按 $x=a\cos\omega t$ 作简谐振动，
假定随机取样时刻在一周期内均匀分布。
求坐标概率密度 $p(x)=\mathrm dP/\mathrm dx$，
其中 $-a<x<a$，并给出图形特征。
\exsolution
对于同一内部位置 $x$，
每周期经过两次，单次停留时间
$\mathrm dt=\mathrm dx/|v_x|$。
因为 $T=2\pi/\omega$，
$v_x^2=\omega^2(a^2-x^2)$，
概率密度
\[
p(x)=\frac{2}{T|v_x|}
=\boxed{\frac1{\pi\sqrt{a^2-x^2}}},\qquad |x|<a.
\]
区间外 $p(x)=0$。
\excheck 归一化
$\int_{-a}^ap(x)\mathrm dx
=\pi^{-1}[\arcsin(x/a)]_{-a}^{a}=1$；
在两端附近密度发散但可积，在 $x=0$ 最小。

\oscproblem{10}{154}
沿同一直线叠加若干同频简谐振动，
用相量法求合振幅：
(a) $x_1=3.0\cos(\omega t+\pi/3)$，
$x_2=8.0\sin(\omega t+\pi/6)$；
(b) $x_1=3.0\cos\omega t$，
$x_2=5.0\cos(\omega t+\pi/4)$，
$x_3=6.0\sin\omega t$。
\exsolution
用复振幅表示 $x=\Re(Ae^{i\omega t})$。
(a) $\sin(\omega t+\pi/6)=\cos(\omega t-\pi/3)$，
所以
$A=3e^{i\pi/3}+8e^{-i\pi/3}$；
相位差 $2\pi/3$，
合振幅
\[
\boxed{|A|=\sqrt{3^2+8^2+2(3)(8)\cos(2\pi/3)}
=\sqrt{49}=7}.
\]
(b) 复振幅
$A=3+5e^{i\pi/4}-6i$。
平方模为
$|A|^2=(3+5/\sqrt2)^2+(5/\sqrt2-6)^2
=70-30/\sqrt2\simeq48.79$，
所以
\[
\boxed{|A|=\sqrt{70-15\sqrt2}\simeq6.99\approx7}.
\]
\excheck (b) 中
$\sin\omega t=\cos(\omega t-\pi/2)$，
因此其相量是 $-6i$ 而不是 $+6i$。

\oscproblem{11}{154}
质点沿同一方向同时参与
$x_1=a\cos\omega t$、
$x_2=a\cos2\omega t$ 两简谐振动。
求合速度的最大模。
\exsolution
总速度
$v=-a\omega[\sin u+2\sin2u]
=-a\omega\sin u(1+4\cos u)$，
$u=\omega t$。
极值条件为
$\mathrm d(\sin u+2\sin2u)/\mathrm du
=\cos u+4\cos2u=0$。
令 $z=\cos u$，
$8z^2+z-4=0$，
有
$z_\pm=(-1\pm\sqrt{129})/16$。
取两候选值及 $u$ 的周期端点，
最大模出现于 $z_+=(-1+\sqrt{129})/16$，
\[
\boxed{v_{\max}=a\omega
\sqrt{1-z_+^2}\,(1+4z_+)
\simeq2.73a\omega}.
\]
\excheck 不能简单将两个分速度振幅
$a\omega$ 与 $2a\omega$ 相加得 $3a\omega$，
因为两个谐波的峰值时刻不同。

\oscproblem{12}{154--155}
质点叠加振动
$x=a\cos(2.1t)\cos(50.0t)$，
$t$ 用秒计。
求两个组成谐波的角频率与拍振周期。
\exsolution
由积化和差：
\[
x=\frac a2[\cos(52.1t)+\cos(47.9t)].
\]
故
\[
\boxed{\omega_1=52.1\,\mathrm{rad/s},
\quad\omega_2=47.9\,\mathrm{rad/s}}.
\]
两频差
$\Delta\omega=4.2\,\mathrm{rad/s}$，
拍振周期
\[
\boxed{T_{\rm beat}=\frac{2\pi}{|\omega_1-\omega_2|}
\simeq1.5\,\mathrm s}.
\]
\excheck 两项等振幅 $a/2$，
包络为 $a|\cos2.1t|$；
相邻包络峰间隔是 $\pi/2.1$ 而非 $2\pi/2.1$。

\oscproblem{13}{155}
质点 A 在参考系 K' 中作某未知频率的简谐振动，
K' 本身沿同一直线相对 K 振动。
当 K' 的频率是 $20$ 或 $24$Hz 时，
在 K 中观察到的拍频均为 $\nu$。
问 K' 振动频率改为多少，拍频变成 $2\nu$？
\exsolution
设 A 相对 K' 的固有振动频率为 $f$。
拍频由两源频率之差给出：
$|f-20|=|f-24|=\nu$，
唯一中点解为
$f=22$Hz、$\nu=2$Hz。
拍频变成 $2\nu=4$Hz 时，
$|f-f'|=4$，故
\[
\boxed{f'=18\,\mathrm{Hz}\ \text{或}\ 26\,\mathrm{Hz}}.
\]
\excheck 两答案都与22Hz差4Hz，
不存在其他非负解。

\oscproblem{14}{155}
平面运动规律
$x=a\sin\omega t$、
$y=b\cos\omega t$，
$a,b,\omega>0$。
求(a)轨迹方程及绕行方向；
(b)加速度矢量关于位置矢量的表达式。
\exsolution
消去时间：
\[
\boxed{\frac{x^2}{a^2}+\frac{y^2}{b^2}=1}.
\]
$t=0$ 时质点在 $(0,b)$，
$v_x=a\omega>0$，
从椭圆顶端向右运动，故沿\emph{顺时针}方向绕行。
分别求二阶导数：
$\ddot x=-\omega^2x$、
$\ddot y=-\omega^2y$，
所以
\[
\boxed{\bm a=-\omega^2\bm r}.
\]
\excheck 即使 $a\ne b$、轨迹为椭圆，
加速度始终指向原点，但速度大小不恒定。

\oscproblem{15}{155}
求平面运动轨迹并作图：
(a) $x=a\sin\omega t$、
$y=a\sin2\omega t$；
(b) $x=a\sin\omega t$、
$y=a\cos2\omega t$。
\exsolution
(a) 令 $u=\omega t$，有
$y=2a\sin u\cos u=2x\cos u$，
而 $\cos^2u=1-x^2/a^2$，
故
\[
\boxed{y^2=4x^2\left(1-\frac{x^2}{a^2}\right)},
\quad |x|\le a.
\]
该轨迹关于两轴对称，是两瓣在原点相交的闭合曲线。
(b) 用
$\cos2u=1-2\sin^2u$：
\[
\boxed{y=a\left(1-\frac{2x^2}{a^2}\right)},
\quad -a\le x\le a.
\]
图形为在有限区间上往复经过的下开口抛物线段。
\excheck (a) 在 $x=0$ 必有 $y=0$，
而 (b) 在 $x=0$ 必有 $y=a$；不可混为同一李萨如图形。

\oscproblem{16}{155}
质量为 $m$ 的质点在一维势场
$U(x)=U_0(1-\cos ax)$ 中，
其中 $U_0,a>0$。
求平衡位置附近的小振动周期。
\exsolution
在稳定平衡点 $x=0$ 附近，
$\cos ax=1-a^2x^2/2+O(x^4)$，
所以
$U(x)\simeq \tfrac12U_0a^2x^2$，
等效弹性系数
$k=U''(0)=U_0a^2$。
角频率
$\omega_0=\sqrt{k/m}=a\sqrt{U_0/m}$，
周期
\[
\boxed{T=\frac{2\pi}{a}\sqrt{\frac m{U_0}}}.
\]
\excheck $a$ 具有长度倒数单位，
所以 $T$ 具有秒的量纲。

\oscproblem{17}{155}
上一题改为
$U(x)=a/x^2-b/x$，其中 $a,b>0$。
求稳定平衡点处小振动周期。
\exsolution
$U'(x)=-2a/x^3+b/x^2$，
稳定点
$x_0=2a/b$。
二阶导数
$U''=6a/x^4-2b/x^3$，
代入
$U''(x_0)=b^4/(8a^3)>0$。
因此
\[
\boxed{T=2\pi\sqrt{\frac{m}{U''(x_0)}}
=\frac{4\pi}{b^2}\sqrt{2ma^3}}.
\]
\excheck 只有正平衡点 $x_0=2a/b$
是本模型内有限稳定点；
二阶导数必须为正才能作小振动。

\oscproblem{18}{155}
一根水平张紧的绳子长
$\ell=1.0\,\mathrm m$，
中点固定一小球，质量 $m=40$g。
绳张力恒定为 $F=10$N。
求小球在竖直平面内的小振动周期。
\exsolution
球竖直位移为 $y$ 时，
每半段绳的长度约为 $\ell/2$，
小偏角
$\sin\theta\simeq y/(\ell/2)=2y/\ell$。
两段绳的回复力
$F_y=-2F\sin\theta\simeq-4Fy/\ell$，
所以
$k_{\rm eff}=4F/\ell$，
\[
\boxed{T=2\pi\sqrt{\frac m{k_{\rm eff}}}
=\pi\sqrt{\frac{m\ell}{F}}
\simeq0.20\,\mathrm s}.
\]
\excheck 原书答案的解析公式与此一致，
但其所印数值“2.0s”与题给40g、10N、1m
直接代入不一致，差十倍；已记录为底本数值疑点。

\oscproblem{19}{155}
绳长 $\ell=20$cm 的单摆，
浸没在密度为摆球密度的 $1/\eta$ 的液体中，
$\eta=3.0$。
忽略流体阻力，求小振动周期。
\exsolution
摆球质量 $m$、体积 $V$，
重力 $mg$，
浮力
$F_b=\rho_\ell Vg=(m/\eta)g$。
有效竖直回复力对应
$g_{\rm eff}=g(1-1/\eta)$，
故
\[
\boxed{T=2\pi\sqrt{
\frac{\ell}{g(1-1/\eta)}}
=2\pi\sqrt{\frac{\eta\ell}{g(\eta-1)}}
\simeq1.1\,\mathrm s}.
\]
\excheck $\eta\to\infty$ 时浮力可忽略，
回到真空中单摆周期；$\eta\downarrow1$ 时周期发散。

\oscproblem{20}{155}
图4.1：长 $\ell$ 的单摆靠近墙壁，静止悬线与墙壁夹小角 $\alpha$；
将摆球再偏离至小角度 $\beta>\alpha$ 后静止释放。
球与墙的碰撞完全弹性。
求其周期。
\exsolution
小角运动满足
$\ddot\theta+\omega_0^2\theta=0$，
$\omega_0=\sqrt{g/\ell}$。
设释放角为 $\theta(0)=\beta$、角速度零，
则
$\theta(t)=\beta\cos\omega_0t$。
按原图的角度约定，碰墙位置是
$\theta=-\alpha$，
到首次碰墙的相位满足
$\cos(\omega_0t_1)=-\alpha/\beta$，
故
$t_1=\omega_0^{-1}[\pi/2+\arcsin(\alpha/\beta)]$。
完全弹性碰撞令角速度反向而不损失机械能，
回到释放端花相同时间，于是
\[
\boxed{T=\frac2{\omega_0}
\left[\frac\pi2+\arcsin\frac\alpha\beta\right]
=2\sqrt{\frac\ell g}
\left[\frac\pi2+\arcsin\frac\alpha\beta\right]}.
\]
\excheck $\alpha\to0$ 时每次仅在摆过竖直位置受墙反弹，周期变为自由单摆周期的一半；
$\alpha\to\beta$ 时碰墙位置靠近另一个转折点，恢复完整简谐周期。
